\documentclass[a4paper,11pt]{article}
\usepackage{amsmath,amssymb,amsthm, tikz,titlesec,hyperref,esint,braket, graphicx, mathrsfs, enumitem}
\usepackage[a4paper,margin=2cm]{geometry}
\linespread{1.3}
\newtheorem{claim}{Claim}[section]

\newcommand\name{Park Chanwoo}   % Name of the student
\newcommand\university{KAIST} % Name of the university
\newcommand\department{Physics} % Name of the department
\newcommand\studentid{20230297} % Student ID
\newcommand\s{\,\;}


\newenvironment{solution}[1]
  {\renewcommand\qedsymbol{$\square$}\begin{proof}[\textbf{Solution#1}]}
  {\end{proof}}
\newenvironment{note}
  {\renewcommand\qedsymbol{$\blacksquare$}\begin{proof}[\textnormal{\textbf{note}}]}
  {\end{proof}}

\title{KAIST\\2026 Special Topics in Physics Quantum Electronic Devices I: Interferometers and Fractional Particles\\
Homework 4\bigskip}
\author{\textbf{\Large \name} \\
% University: \university\\
Department: \department\\
Student ID: \studentid}
\date{\today}

\begin{document}
\thispagestyle{empty}
\maketitle
\tableofcontents
\titleformat{\section}[frame]{\pagebreak}{\filright
\footnotesize  \enspace \textsf{KAIST --- Interferometers and Fractional Particles 2026 Spring}\enspace}{6pt}{\Large\bfseries\filcenter}
\setcounter{secnumdepth}{1}
\setcounter{tocdepth}{2}

\newcommand{\tr}{\operatorname{tr}}
\newcommand{\sgn}{\operatorname{sgn}}
\newcommand{\supp}{\operatorname{supp}}
\renewcommand{\Re}{\operatorname{Re}}
\renewcommand{\Im}{\operatorname{Im}}
\newcommand{\boltz}{k_{\mathrm{B}}}

\section{\#1}
\subsection{(1a)}

From the unitary condition for the scattering matrix, we can find the following relations:
\begin{align}
    |t|^2 + |r|^2 &= 1, \\
    t^*r + r^*t &= 0.
\end{align}

Since the state of bosonic particles must be symmetrized, The incoming two-particle state can be written as
\begin{equation}
    \ket{\text{incoming}} = \frac{1}{\sqrt{2}}\left(\ket{\psi_h}\ket{\psi_v} + \ket{\psi_v}\ket{\psi_h}\right).
\end{equation}

After the beam splitter, the state becomes
\begin{align}
    \ket{\text{outgoing}} 
    &= S\otimes S \ket{\text{incoming}} \\
    &= \frac{1}{\sqrt{2}}\left(S\ket{\psi_v}S\ket{\psi_h} + S\ket{\psi_h}S\ket{\psi_v}\right) \\
    &= \frac{1}{\sqrt{2}}\left(t\ket{\psi_v} + r\ket{\psi_h}\right)\left(r\ket{\psi_v} + t\ket{\psi_h}\right) + \frac{1}{\sqrt{2}}\left(r\ket{\psi_v} + t\ket{\psi_h}\right)\left(t\ket{\psi_v} + r\ket{\psi_h}\right) \\
    &= \frac{t^2 + r^2}{\sqrt{2}}\left(\ket{\psi_v}\ket{\psi_h} + \ket{\psi_h}\ket{\psi_v}\right) + \frac{2tr}{\sqrt{2}}\left(\ket{\psi_v}\ket{\psi_v} + \ket{\psi_h}\ket{\psi_h}\right)
\end{align}

Hence, the probabilities are:
\begin{align}
    P(2, 0) &= 2RT, \\
    P(0, 2) &= 2RT, \\
    P(1, 1) &= (t^2 + r^2)({t^*}^2 + {r^*}^2) \\
    &= R^2 + T^2 + r^2{t^*}^2 + t^2{r^*}^2 \\
    &= R^2 + T^2 - 2RT + (r^*t + t^*r)^2 \\
    &= (R - T)^2.
\end{align}

\subsection{(1b)}

For fermionic particles, antisymmetrizing the state,
\begin{align}
    \ket{\text{outgoing}} 
    &= S\otimes S \ket{\text{incoming}} \\
    &= \frac{1}{\sqrt{2}}\left(S\ket{\psi_v}S\ket{\psi_h} - S\ket{\psi_h}S\ket{\psi_v}\right) \\
    &= \frac{1}{\sqrt{2}}\left(t\ket{\psi_v} + r\ket{\psi_h}\right)\left(r\ket{\psi_v} + t\ket{\psi_h}\right) - \frac{1}{\sqrt{2}}\left(r\ket{\psi_v} + t\ket{\psi_h}\right)\left(t\ket{\psi_v} + r\ket{\psi_h}\right) \\
    &= \frac{t^2 - r^2}{\sqrt{2}}\left(\ket{\psi_v}\ket{\psi_h} - \ket{\psi_h}\ket{\psi_v}\right)
\end{align}
Hence, the probabilities are:
\begin{align}
    P(2, 0) &= 0, \\
    P(0, 2) &= 0, \\
    P(1, 1) &= (t^2 - r^2)({t^*}^2 - {r^*}^2) \\
    &= R^2 + T^2 - r^2{t^*}^2 - t^2{r^*}^2 \\
    &= R^2 + T^2 + 2RT - (r^*t + t^*r)^2 \\
    &= (R + T)^2.
\end{align}


\subsection{(1c)}

\subsection{(1c)}

Using second quantization, the problem gives
\begin{eqnarray}
    \begin{pmatrix}
        d_1^\dagger \\
        d_2^\dagger
    \end{pmatrix}
    =
    \begin{pmatrix}
        r & t \\
        t & r
    \end{pmatrix}
    \begin{pmatrix}
        d_3^\dagger \\
        d_4^\dagger
    \end{pmatrix}
\end{eqnarray}

i.e., $d_1^\dagger = r d_3^\dagger + t d_4^\dagger$ and $d_2^\dagger = t d_3^\dagger + r d_4^\dagger$.
The incoming state transforms as
\begin{align}
    d_1^\dagger d_2^\dagger \ket{0}
    &= (rd_3^\dagger + td_4^\dagger)(td_3^\dagger + rd_4^\dagger)\ket{0} \\
    &= rt (d_3^\dagger)^2 + r^2 d_3^\dagger d_4^\dagger + t^2 d_4^\dagger d_3^\dagger + tr (d_4^\dagger)^2
\end{align}


For the bosonic commutation relation, $[d_i^\dagger, d_j^\dagger] = 0$,
\begin{align}
    d_1^\dagger d_2^\dagger\ket{0} = (t^2 + r^2) d_3^\dagger d_4^\dagger\ket{0} + \sqrt{2}tr \frac{(d_3^\dagger)^2}{\sqrt2}\ket{0} + \sqrt{2}tr \frac{(d_4^\dagger)^2}{\sqrt2}\ket{0}    
\end{align}
gives probabilities of $P(2, 0) = P(0, 2) = 2RT$ and $P(1, 1) = (R - T)^2$.

For the fermionic anticommutation relation, $\{d_i^\dagger, d_j^\dagger\} = 0$,
\begin{align}
    d_1^\dagger d_2^\dagger\ket{0} = (t^2 - r^2) d_3^\dagger d_4^\dagger\ket{0}
\end{align}
gives probabilities of $P(2, 0) = P(0, 2) = 0$ and $P(1, 1) = 1$.

The results are consistent with the first quantization approach.

\end{document}
